\documentclass[14pt,a4paper,oneside]{extbook}
\usepackage[utf8]{vietnam}
\usepackage{amsmath,amssymb,mathrsfs,fancyhdr,enumerate,multirow,makecell}
\usepackage{tkz-euclide,tikz-3dplot,tikz,tkz-tab}
\usetikzlibrary{arrows.meta,calc,intersections,angles,quotes,backgrounds,shapes.geometric,patterns,through,shadows,positioning}
\usepackage{yhmath}
\usepackage[top=0.7cm, bottom=0.7cm, left=1.0cm, right=1.0cm, headheight=16pt, headsep=3mm, footskip=5mm]{geometry}
\usepackage[hidelinks,unicode]{hyperref}
\usepackage{qrcode}
\usepackage[many]{tcolorbox}
\usepackage{fontawesome}

\pagestyle{fancy}
\fancyhf{}
\fancyhead[L]{\footnotesize\bfseries\color{teal!85!black}\faBook~SỔ TAY KIẾN THỨC TOÁN 5}
\fancyhead[R]{\footnotesize\bfseries\color{blue!85!black}Lớp Toán Thầy Lĩnh -- loptoanthaylinh.vn}
\fancyfoot[C]{\footnotesize\bfseries\color{gray!80!black}-- Sổ tay kiến thức: Bài 67. Luyện tập chung thống kê, xác suất --}
\renewcommand{\headrulewidth}{0.4pt}

\begin{document}

% --- BANNER TIÊU ĐỀ BÀI HỌC ---
\noindent
\begin{tcolorbox}[
  enhanced,
  colback=teal!8!white,
  colframe=teal!80!black,
  arc=3mm,
  boxrule=1.1pt,
  width=\linewidth,
  halign=center,
  top=2.2mm, bottom=2.2mm,
  left=3mm, right=3mm,
  drop fuzzy shadow
]
  {\large\bfseries\color{teal!95!black}\faBookmark\ BÀI 67. LUYỆN TẬP CHUNG THỐNG KÊ \& XÁC SUẤT}\\[2pt]
  {\scriptsize\color{black!75}\itshape Tổng ôn các dạng bài toán về bảng số liệu, biểu đồ quạt tròn và xác suất thực nghiệm trên 1 trang A4}
\end{tcolorbox}

\vspace{3.2mm}
% --- KHỐI 1: HAI TRỤ CỘT CHỦ ĐỀ 11 ---
\noindent
\begin{minipage}[t]{0.54\textwidth}
  \begin{tcolorbox}[
    enhanced,
    title={\footnotesize\bfseries\color{white}\faBarChart\ 1. Trọng tâm Thống kê},
    colback=blue!2!white,
    colframe=blue!75!black,
    arc=2.5mm, boxrule=0.8pt, top=1.8mm, bottom=1.8mm, left=2.2mm, right=2.2mm
  ]
    \footnotesize
    $\bullet$ \textbf{Bảng số liệu:} Đọc đúng dòng, cột và tìm số liệu nhiều nhất, ít nhất.\\
    $\bullet$ \textbf{Biểu đồ quạt tròn:} Chuyển đổi linh hoạt giữa tỉ số phần trăm và số lượng thành phần ($A = T \times p\%$).
  \end{tcolorbox}
\end{minipage}%
\hfill
\begin{minipage}[t]{0.44\textwidth}
  \begin{tcolorbox}[
    enhanced,
    title={\footnotesize\bfseries\color{white}\faCubes\ 2. Trọng tâm Xác suất},
    colback=teal!2!white,
    colframe=teal!75!black,
    arc=2.5mm, boxrule=0.8pt, top=1.8mm, bottom=1.8mm, left=2mm, right=2mm
  ]
    \footnotesize
    $\bullet$ \textbf{Xác suất thực nghiệm:} Tỉ số $=\frac{m}{n}$ với $m$ là số lần sự kiện xảy ra, $n$ là tổng số lần thử.\\
    $\bullet$ Dự đoán số lần xuất hiện khi thực hiện thí nghiệm ở quy mô lớn hơn.
  \end{tcolorbox}
\end{minipage}

\vspace{3mm}

% --- KHỐI 3: SƠ ĐỒ TƯ DUY ---
\noindent
\begin{tcolorbox}[
  enhanced,
  title={\footnotesize\bfseries\color{white}\faSitemap\ Sơ đồ tư duy: Hệ thống hóa Chủ đề 11 -- Thống kê \& Xác suất},
  colback=yellow!2!white,
  colframe=teal!85!black,
  arc=3mm, boxrule=0.9pt, top=1.8mm, bottom=1.8mm, left=1.5mm, right=1.5mm,
  halign=center
]
  \begin{tikzpicture}[
    scale=0.92,
    every node/.style={transform shape},
    root/.style={rectangle, rounded corners=5pt, fill=teal!85!black, text=white, font=\bfseries\footnotesize, inner sep=5pt, align=center},
    b1/.style={rectangle, rounded corners=4pt, fill=blue!12, draw=blue!75!black, text=blue!90!black, font=\bfseries\scriptsize, inner sep=4.5pt, align=center},
    b2/.style={rectangle, rounded corners=4pt, fill=green!15, draw=green!65!black, text=green!80!black, font=\bfseries\scriptsize, inner sep=4.5pt, align=center}
  ]
    \node[root] (core) at (0, 0.85) {\faBarChart\ THỐNG KÊ \& XÁC SUẤT};
    \node[b1] (n1) at (-4.8, -0.45) {\textbf{PHẦN 1: THỐNG KÊ}\\[2pt]Thu thập $\rightarrow$ Phân loại $\rightarrow$ Lập bảng\\Biểu đồ quạt: $A = T \times p\%$};
    \node[b2] (n2) at (4.8, -0.45) {\textbf{PHẦN 2: XÁC SUẤT}\\[2pt]Tỉ số lần lặp lại sự kiện: $\frac{m}{n}$\\Biểu diễn dạng phân số, STP, phần trăm};

    \draw[teal!85!black, thick, -{Stealth[length=2mm]}] (core) -| (n1);
    \draw[teal!85!black, thick, -{Stealth[length=2mm]}] (core) -| (n2);
  \end{tikzpicture}
\end{tcolorbox}

\vspace{3mm}

% --- KHỐI 4: BÀI TOÁN DỰ ĐOÁN & BẪY TRẮC NGHIỆM ---
\noindent
\begin{minipage}[t]{0.54\textwidth}
  \begin{tcolorbox}[
    enhanced,
    title={\footnotesize\bfseries\color{white}\faLineChart\ Dạng toán dự đoán số lần xuất hiện},
    colback=green!3!white,
    colframe=green!60!black,
    arc=2.5mm, boxrule=0.8pt, top=1.8mm, bottom=1.8mm, left=2.2mm, right=2.2mm
  ]
    \footnotesize
    Nếu tỉ số xuất hiện của sự kiện là $p\%$, khi thực hiện $N$ lần thử:\\
    $\bullet$ Số lần dự đoán xuất hiện: $\mathbf{k \approx N \times p\%}$.\\
    \textit{VD:} Tung bóng $50$ lần trúng $20$ lần ($40\%$). Dự đoán tung $100$ lần trúng khoảng: $100 \times 40\% = 40$ lần.
  \end{tcolorbox}
\end{minipage}%
\hfill
\begin{minipage}[t]{0.44\textwidth}
  \begin{tcolorbox}[
    enhanced,
    title={\footnotesize\bfseries\color{white}\faExclamationTriangle\ Mẹo tránh bẫy trắc nghiệm},
    colback=red!2!white,
    colframe=red!70!black,
    arc=2.5mm, boxrule=0.8pt, top=1.8mm, bottom=1.8mm, left=2.2mm, right=2.2mm
  ]
    \footnotesize
    $\bullet$ Xác suất thực nghiệm mang tính \textbf{xấp xỉ}, không phải kết quả tuyệt đối chắc chắn.\\
    $\bullet$ Đọc kĩ đề hỏi số lần trúng hay số lần trượt!
  \end{tcolorbox}
\end{minipage}

\end{document}
